El artículo \textit{Identifying group contributions in NBA lineups with spectral analysis}, desarrollado por Devlin y Uminsky (2020), se centra en analizar cómo las diferentes combinaciones de jugadores contribuyen al rendimiento de un equipo de baloncesto. Los autores plantean que las estadísticas individuales no siempre representan el impacto real que tiene un jugador dentro de la cancha, ya que su rendimiento también depende de los compañeros con los que comparte una alineación.

El estudio utiliza datos \textit{play-by-play} de la temporada 2015--2016 de la NBA, tomando como referencia los 15 jugadores de cada equipo con mayor número de posesiones disputadas. A partir de estos datos calculan el indicador \textit{plus-minus}, que representa el diferencial entre los puntos anotados y recibidos por una alineación durante su participación. Considerando una plantilla de 15 jugadores, existen 3003 combinaciones posibles de cinco jugadores, aunque muchas de ellas no llegan a presentarse durante los partidos.

Para analizar estas combinaciones, los autores hacen uso del análisis espectral, una técnica matemática que permite separar las contribuciones al rendimiento del equipo dependiendo del número de jugadores que participan en una combinación. De esta manera, el rendimiento se divide en diferentes niveles:

\begin{itemize}
    \item Contribución individual de cada jugador.
    \item Contribución de parejas de jugadores.
    \item Contribución de grupos de tres jugadores.
    \item Contribución de grupos de cuatro jugadores.
    \item Contribución de alineaciones completas de cinco jugadores.
\end{itemize}

Como principales ventajas de este método es que permite identificar si el buen rendimiento de una combinación se debe a la calidad individual de sus integrantes o a la interacción que existe entre ellos. Por ejemplo, una pareja puede presentar buenos resultados porque uno de sus jugadores tiene un rendimiento individual elevado, sin que necesariamente exista una contribución positiva adicional por jugar juntos. 

En los resultados se analizaron las contribuciones de diferentes grupos de jugadores de los Golden State Warriors y los Boston Celtics durante la temporada estudiada. Los autores muestran que algunas combinaciones presentan un rendimiento diferente al que se podría esperar observando únicamente sus estadísticas individuales. Permitiendo identificar grupos con interacciones favorables y también combinaciones cuyo rendimiento se explica principalmente por las capacidades individuales de sus integrantes.

Como ventaja, este estudio permite evaluar la interacción entre jugadores sin limitarse únicamente a las alineaciones completas de cinco integrantes. Además, proporciona información para analizar posibles rotaciones, ya que permite identificar parejas o grupos que presentan buenas contribuciones colectivas. Otra característica importante es que no necesita clasificar previamente a los jugadores según su posición o estilo de juego, debido a que las contribuciones se obtienen directamente del rendimiento observado de las alineaciones.

Sin embargo, una de sus limitaciones es que el análisis se realiza a partir de datos históricos y del rendimiento que tuvieron las combinaciones durante los partidos. 

Otra limitación es que las combinaciones de jugadores no participan durante la misma cantidad de posesiones. Algunas alineaciones tienen una gran cantidad de información disponible, mientras que otras aparecen pocas veces, lo que dificulta la comparación de sus contribuciones. Aunque los autores utilizan medidas que consideran el número de posesiones para facilitar el análisis, esta diferencia continúa siendo una característica importante de los datos utilizados.

\cite{devlin2020identifyinggroupcontributionsnba}